\section{Que molde estatístico descreve cada indicador}
\label{sec:distribuicoes}

Cada indicador tem uma forma própria de variar entre hospitais, e há
moldes matemáticos (as chamadas famílias de distribuição) desenhados
para cada tipo de forma. Identificar o molde certo é importante porque
toda análise posterior fica mais confiável quando respeita a natureza
do dado: proporções entre 0 e 1 pedem um molde, valores em reais com
cauda longa pedem outro, contagens pedem um terceiro. Esta seção testa
os moldes contra os dados e documenta as escolhas.

\subsection{Mortalidade: uma proporção com dois grupos de hospitais}

A mortalidade é uma proporção (fica entre 0 e 1), e o molde clássico
para proporções é a distribuição \textbf{Beta}. O ajuste de uma Beta
única descreve bem a escala geral, mas erra o pico de hospitais com
mortalidade quase nula. A Figura~\ref{fig:mistura} mostra a solução e,
ao mesmo tempo, a explicação: somando \textbf{duas Betas} (uma para
cada grupo de hospitais), a curva acompanha o histograma com nitidez, e
a nota de ajuste melhora de forma expressiva (ganho de 1.124 pontos de
AIC). O ponto central para a gestão é que esses dois grupos
\textbf{existem de verdade no painel}: um grupo de \textbf{58 hospitais
de baixo óbito estrutural} (mortalidade mediana abaixo de 2\%,
majoritariamente maternidades, pediátricos e oncológicos) e os
\textbf{256 demais}, hospitais gerais com mortalidade em torno de 6\%.
A aparente anomalia do histograma é, portanto, apenas o retrato de dois
perfis assistenciais diferentes convivendo na rede. Como também existem
observações com mortalidade exatamente zero (1,5\%), o molde completo
para a mortalidade é a \textbf{Beta com um tratamento especial para os
zeros}.

\begin{figure}[H]
\centering
\subfig{fig_ae_12_mistura_mort_all.png}{Um molde só vs a soma de dois}
\subfig{fig_ae_12_decomposicao_mort_all.png}{Os dois grupos reais de
hospitais}
\caption{Mortalidade geral: à esquerda, a curva cinza tracejada é o
molde único (que erra o pico perto de zero) e a curva escura é a soma
de dois moldes, que acompanha o histograma; à direita, a mesma soma
reconstruída a partir dos dois grupos reais de hospitais do painel.}
\label{fig:mistura}
\end{figure}

\comoler{no painel da esquerda, compare as duas curvas com as barras: a
tracejada passa por cima do vale e erra os dois picos; a escura desenha
os dois picos. No painel da direita, as curvas tracejadas coloridas são
os dois grupos de hospitais (baixo óbito à esquerda, demais à direita)
e a curva preta é a soma dos dois, praticamente idêntica ao histograma:
a forma estranha da mortalidade é só a mistura de dois perfis de
hospital.}

\subsection{Alta complexidade: o peso dos zeros}

A fração de alta complexidade também é uma proporção, mas com uma
particularidade dominante: \textbf{um terço das observações vale
exatamente zero} (hospitais que não realizam alta complexidade) e uma
pequena parcela vale exatamente um (hospitais que só realizam alta
complexidade). O molde adequado trata três situações em separado: a
chance de o hospital estar no zero, a chance de estar no um, e a
distribuição de quem está no meio; esse molde combinado é conhecido
pela sigla \textbf{ZOIB}. A Figura~\ref{fig:betapct} mostra o ajuste do
miolo e a Figura~\ref{fig:zeros} mostra que \textbf{os zeros não são
aleatórios}: a proporção de hospitais sem alta complexidade cai de
82,6\% na faixa 2 de complexidade para zero na faixa 6, e de 41,1\% nos
Filantrópicos para 10,3\% nas OSS. Ou seja, saber a complexidade, o
porte e a categoria do hospital já diz muito sobre a chance de ele
estar no grupo do zero.

\begin{figure}[H]
\centering
\includegraphics[width=.96\textwidth]{fig_ae_02_beta_pct_alta_complex.png}
\caption{Fração de alta complexidade entre os hospitais que realizam
alguma alta complexidade: histograma e molde Beta ajustado. O eixo
vertical está em escala logarítmica para a cauda inteira ficar
visível.}
\label{fig:betapct}
\end{figure}

\comoler{a grande maioria dos hospitais que fazem alta complexidade faz
pouca (a concentração à esquerda; metade fica abaixo de 3,7\%), e uma
minoria é altamente especializada (as barras à direita, perto de
100\%). A curva escura é o molde Beta, que acompanha essa queda
contínua.}

\begin{figure}[H]
\centering
\subfig{fig_ae_15_zeros_pct_alta_categoria.png}{Por categoria de
gestão}
\subfig{fig_ae_15_zeros_pct_alta_faixa.png}{Por faixa de complexidade}
\caption{Percentual de observações com alta complexidade exatamente
zero, por grupo de hospitais.}
\label{fig:zeros}
\end{figure}

\comoler{cada barra é a proporção de hospitais ano daquele grupo que
não realizou nenhuma internação de alta complexidade. A queda das
barras da esquerda para a direita no painel (b) mostra que o zero é
questão de perfil e estrutura: quase todos os hospitais de faixa 2
estão no zero, e nenhum de faixa 6 está.}

\subsection{Faturamento e tempo de permanência: valores positivos com cauda
longa}

Faturamento por saída e TMP são valores positivos com cauda à direita. Cinco
moldes clássicos foram postos à prova
(Tabela~\ref{tab:familias})\footnote{Guia dos tamanhos de amostra que
aparecem no documento: \textbf{3.454} é o painel completo;
\textbf{3.453} exclui a única observação com TMP igual a zero (Hospital
Infante D.\ Henrique, 2021), incompatível com a escala logarítmica;
\textbf{3.255} exclui também as 198 observações dos 18 hospitais de
longa permanência; \textbf{3.452} exclui as duas observações de
denominador frágil da Seção~\ref{sec:qualidade}, regra usada nos
modelos da fase de estimação; e \textbf{2.571} são as observações com
UTI ativa depois de retiradas essas duas frágeis. Cada conta foi
conferida no painel.},
comparados pela nota AIC e pelo gráfico QQ. O resultado é consistente:
o molde \textbf{log logístico (também chamado de Fisk)} vence nas três
amostras, com a \textbf{LogNormal} logo atrás, e a Gama, molde usado
por tradição em análises de custo, fica sempre entre as piores. No TMP
há um detalhe decisivo: enquanto os hospitais de longa permanência
estiverem na amostra, nenhum molde acompanha o teto de 30,5 dias;
\textbf{retirando os 18 hospitais de longa permanência, o ajuste do
molde log logístico fica praticamente perfeito} (r de 0,9975 na
Figura~\ref{fig:qqtmpsemlp}). A leitura prática: faturamento e TMP devem ser
analisados na escala logarítmica, e os hospitais de longa permanência
merecem tratamento em separado nas análises de TMP. Os gráficos QQ
completos estão no Apêndice~\ref{ap:diagnosticos}: no faturamento, o molde
log logístico acompanha os dados por quase toda a extensão (r de
0,9754, Figura~\ref{fig:qqcusto}); no TMP do painel completo, nenhum
molde reproduz o teto de 30,5 dias (Figura~\ref{fig:qqtmp}).

\begin{table}[H]
\centering
\caption{Comparação dos cinco moldes de distribuição para faturamento e TMP
pela nota AIC (Seção~\ref{sec:painel}) e pelo r do QQ
(Apêndice~\ref{ap:diagnosticos}).}
\label{tab:familias}
\small
\begin{tabular}{llrrr}
\toprule
Amostra & Molde & $n$ & AIC & r do QQ \\
\midrule
Faturamento por saída & Log logística (Fisk) & 3.454 & 55.483 & 0,9754 \\
                & LogNormal            & 3.454 & 55.547 & 0,9162 \\
                & Gaussiana inversa    & 3.454 & 55.608 & 0,9018 \\
                & Gama                 & 3.454 & 56.499 & 0,8411 \\
                & Weibull              & 3.454 & 56.995 & 0,8415 \\
\midrule
TMP, painel completo & Log logística (Fisk) & 3.453 & 15.990 & 0,8597 \\
                     & LogNormal            & 3.453 & 16.650 & 0,8696 \\
                     & Gaussiana inversa    & 3.453 & 16.767 & 0,8736 \\
                     & Gama                 & 3.453 & 17.990 & 0,8379 \\
                     & Weibull              & 3.453 & 18.701 & 0,8399 \\
\midrule
TMP sem longa permanência & Log logística (Fisk) & 3.255 & 12.147 & 0,9975 \\
                          & LogNormal            & 3.255 & 12.226 & 0,9810 \\
                          & Gaussiana inversa    & 3.255 & 12.260 & 0,9787 \\
                          & Gama                 & 3.255 & 12.417 & 0,9635 \\
                          & Weibull              & 3.255 & 13.204 & 0,9403 \\
\bottomrule
\end{tabular}
\end{table}

\begin{figure}[H]
\centering
\subfig{fig_ae_03_hist_log_custo_saida.png}{Faturamento por saída}
\subfig{fig_ae_03_hist_log_tmp.png}{TMP}
\caption{Distribuições de faturamento e TMP com o eixo horizontal em escala
logarítmica (os rótulos mostram os valores reais).}
\label{fig:histlog}
\end{figure}

\comoler{na escala logarítmica, o faturamento (a) vira um morro quase
simétrico, sinal clássico de que a análise deve ser feita nessa escala.
No TMP (b), além do morro principal entre 3 e 8 dias, aparece um
segundo morrinho isolado à direita, junto de 30 dias: é a subpopulação
de longa permanência.}

\begin{figure}[H]
\centering
\includegraphics[width=.96\textwidth]{fig_ae_03_qq_fisk_tmp_sem_longa_perm.png}
\caption{Gráfico QQ do TMP sem os 18 hospitais de longa permanência: os
pontos ficam praticamente sobre a reta (r de 0,9975).}
\label{fig:qqtmpsemlp}
\end{figure}

\comoler{este é o gráfico QQ do TMP refeito sem os 18 hospitais de
longa permanência (a versão com o painel completo está na
Figura~\ref{fig:qqtmp}, no apêndice): o achatamento desaparece e o
molde descreve os
dados quase perfeitamente. É a demonstração visual de que o problema
não era o molde, era a mistura de dois perfis de hospital.}

\subsection{Produção: contagens muito mais variáveis que o usual}

A produção anual (número de saídas) é uma contagem. O molde mais
simples para contagens, a Poisson, supõe que a variação entre hospitais
seja comparável à média; aqui \textbf{a variação é milhares de vezes
maior que a média} (razão de 6.381 no agregado e de 93 mesmo comparando
cada hospital só consigo mesmo ao longo dos anos). O molde adequado é a
\textbf{Binomial Negativa}, que aceita essa variação extra.

\subsection{Os moldes melhoram quando se compara o comparável}

Uma pergunta natural: se os QQ do painel completo não são perfeitos, o
molde está errado? As Figuras~\ref{fig:qqcondA} a~\ref{fig:qqcondC},
reunidas no Apêndice~\ref{ap:diagnosticos},
respondem que não: o desajuste vem de misturar 314 hospitais muito
diferentes num único gráfico. Quando se compara cada hospital consigo mesmo (controlando o
efeito do hospital e do ano), os desvios praticamente desaparecem: no
faturamento, a assimetria restante é quase nula; na mortalidade, a forte
assimetria da distribuição bruta na escala \textbf{logito} (a escala
interna em que o modelo trabalha proporções), de 1,72 negativa, cai
para 0,11 depois do controle; na escala original da
Tabela~\ref{tab:univariada} a assimetria é outra (0,72 positiva), e os
dois números não se contradizem. Em resumo, \textbf{os moldes escolhidos descrevem bem os
dados quando se compara o comparável}.
